\begin{abstract}
This project investigates whether the shape of a dielectric resonator could be
optimised to maximise the Purcell factor of a solid-state MASER at room
temperature and zero field. COMSOL Multiphysics\textregistered{}{} and a
Java-based genetic algorithm (GA) were used to simulate and test thousands of
resonator shapes to determine those that optimised electromagnetic confinement
while still being resonant at 1.45 GHz. A 10x10 binary subgrid design using a
binary matrix allowed for shape exploration without remeshing. Results showed
that in theoretical, loss-free conditions, the optimised geometry significantly
outperformed the original MASER-in-a-Shoebox design in Purcell factor and
simulated masing output power. However, when the dielectric is assigned a
realistic loss tangent, this outperformance was not seen and the GA could not
find dielectric shapes which operate at the correct mode. Additional future work
would require exploring transfer to COMSOL LiveLink\texttrademark{} for
MATLAB\textregistered{}{}, which would result in more robust algorithms and
co-localisation with machine learning methods. Overall, this work reassures that
resonator geometry has a strong influence on MASER performance and provides a
foundation for more advanced, physically realistic optimisation of cavity
geometries for future microwave technologies.
\end{abstract}
